\chapter{The pipeline: from files to figures}\label{ch:pipeline}

\begin{idea}
Ten numbered scripts run in order; each reads the Parquet files made by earlier ones and writes either more Parquet or a table and figure into \file{results/}. If you understand what each step consumes and produces, you can change one without breaking the rest.
\end{idea}

\section{Tools and why}
\begin{description}
  \item[DuckDB] runs SQL directly on Parquet files without loading them into memory; it is what makes joins over 400 million votes feasible on a laptop \src{Raasveldt \& M\"uhleisen 2019}.
  \item[pandas / NumPy / SciPy / statsmodels] for the statistics on the small tables DuckDB returns.
  \item[lifelines] Kaplan--Meier and Cox models \src{Davidson-Pilon 2019}.
  \item[matplotlib] figures, with shared styling in \file{plotstyle.py}.
  \item[uv] manages the Python environment from \file{code/pyproject.toml} and \file{uv.lock}.
\end{description}

\section{The scripts}
All scripts are in \file{code/experiments/01-reply-opposition-lifetime/}. \file{common.py} holds paths, the censoring constants, a DuckDB connection with memory settings, and the \texttt{log()} helper that writes to \file{results/run.log}.

\begin{longtable}{@{}p{3.2cm} p{4.2cm} p{3.6cm} p{3.4cm}@{}}
\toprule
Script & Does & Reads & Writes \\
\midrule
\endhead
\file{00\_to\_parquet} & Converts raw TSVs to typed Parquet; hex ids to integers; NA to NULL. Skips outputs that exist. & raw tables & \file{derived/*.parquet} \\
\file{01\_reply\_edges} & Builds reply edges A$\to$B; applies the exclusions of Chapter~\ref{ch:data}. & comments & \file{reply\_edges}; funnel CSV \\
\file{02\_null\_pairs} & One null target per edge (hash-based, same article, earlier comment). & reply edges, comments & null pairs \\
\file{03\_dyads} & Vote totals for each pair, in both directions, for real and null pairs; A's vote on the parent. & votes, edges, null & \file{dyads}, \file{edges\_stance} \\
\file{04\_h1\_comment\_level} & H1 by Mantel--Haenszel, by year, depth, topic. & edges, votes & \file{04\_*.csv/png} \\
\file{05\_h2\_h3\_dyad\_faction} & H2 (Beta edges, null, bootstrap), H3 (factions), reply-back, projection. & dyads, stance, factions & \file{05\_*.csv/png} \\
\file{06\_lifetime} (+ \file{h4.py}) & Lifetime, KM, early features, Cox. & comments, votes, edges & \file{06\_*.csv/png} \\
\file{07\_multiverse} & 24 H2 and 18 H4 alternative specifications. & as 05, 06 & \file{07\_*} \\
\file{08\_positive\_votes} & Mirror of H1--H4 for upvotes. & as above & \file{08\_*} \\
\file{09\_top\_decile\_lifetime} & Long-lived group analysis. & H4 population & \file{09\_*} \\
\file{check\_duplicate\_comments}, \file{check\_self\_votes} & Sanity checks that are not part of the chain. & raw & console \\
\bottomrule
\end{longtable}

\section{How to run it}
From the \file{code/} folder (Windows, PowerShell or Git Bash):
\begin{lstlisting}
cd code
uv sync                                   # create/refresh the environment
cd experiments/01-reply-opposition-lifetime
uv run python 00_to_parquet.py            # slow, one time
uv run python 01_reply_edges.py
uv run python 02_null_pairs.py
uv run python 03_dyads.py
uv run python 04_h1_comment_level.py
uv run python 05_h2_h3_dyad_faction.py
uv run python 06_lifetime.py
uv run python 07_multiverse.py
uv run python 08_positive_votes.py
uv run python 09_top_decile_lifetime.py
\end{lstlisting}
Each script prints timed log lines and appends them to \file{results/run.log}. The parquet conversion and the dyad aggregation are the slow steps; the analysis scripts take minutes. The memory limit in \file{common.py} is 22\,GB; lower it if your machine has less RAM.

\begin{pitfall}
The derived Parquet files are not versioned. If you change an early script (for example the cleaning) you must rerun \emph{everything downstream}, or a later script will silently use stale data. A sensible habit: after any change to \file{01}--\file{03}, rerun \file{04} onward and compare with the previous \file{results/} (\texttt{git} makes this easy once you commit).
\end{pitfall}

\section{Reading a script}
A typical script has three parts. (1) A docstring that states the question; read it first, it is the best summary of the script. (2) DuckDB SQL that builds one row per unit of analysis (a comment, a pair, a user). (3) Python that turns the table into estimates, intervals and a figure. When a number in the thesis looks wrong, find the SQL that produced the table and print its first rows before suspecting the statistics.

\begin{tryit}
Open \file{04\_h1\_comment\_level.py}. Find (a) the query that defines the strata, (b) the function that computes $\widehat{OR}_{MH}$, (c) the bootstrap loop. Check that (b) matches the formula in Section~\ref{sec:h1}.
\end{tryit}

\section{The documents around the code}
\begin{tabularx}{\textwidth}{@{}l Y@{}}
\toprule
File (in \file{notes/experiments/01-reply-opposition-lifetime/}) & Use \\
\midrule
\file{PLAN.md} & the original plan and the decisions it made \\
\file{NOTES.md} & numbered decisions with their rationale and status tags \\
\file{RESULTS.md} & every number, by section; the source for the thesis tables \\
\file{EXPLICACION.md}, \file{EXPLICACION-RESUMEN.md} & Spanish explanations of the method and a short version \\
\file{PREGUNTAS.md} & open questions \\
\file{CODE-REVIEW.md} & the line-by-line review (Chapter~\ref{ch:critique}) \\
\file{results/} & CSVs, figures, \file{run.log} \\
\bottomrule
\end{tabularx}

\section{The eighteen choices}\label{sec:choices}
Every decision that someone could challenge in a defence, with where it is justified. Use this list to prepare answers.

\begin{longtable}{@{}l p{4.7cm} p{6.4cm}@{}}
\toprule
 & Choice & Reason (short) \\
\midrule
\endhead
C1 & Votes as stance proxy & No text; votes are the only signal of approval \\
C2 & Reply = edge A$\to$B via parent id & Observed, not inferred \\
C3 & Exclude deleted authors and self-replies & Cannot attribute; self-replies carry no opposition \\
C4 & Votes cleaned (deleted voters/authors, duplicates) & Needed to join; the dataset authors recommend it \\
C5 & Co-present null, article level & Same article and exposure; only choice changes \\
C6 & Stratify H1 by comment & Comparison on the same comment \\
C7 & Beta(1,2) prior, thresholds 0.4/0.6 & Inherited from FAULTANA \\
C8 & $\ge5$ votes for a defined edge & One vote should not define a relation \\
C9 & ``Excl'' variant: drop vote on replied comment & So the reply does not explain its own vote \\
C10 & Bootstrap by replier, Poisson weights & Edges of one user are dependent \\
C11 & Lifetime and censoring rules & Window ends 2022-12; see Chapter~\ref{ch:data} \\
C12 & Landmark at day 30 & Avoid immortal-time bias \\
C13 & Cox stratified by cohort & Different baselines \\
C14 & Shrinkage with $M=2$ & Small samples; fixed an undefined share \\
C15 & Multiverse of specifications & Show fragility \\
C16 & Effect sizes and CIs, no $p$-values & Millions of rows \\
C17 & Top-10\% group definitions \emph{[AI, not yet reviewed]} & See Section~\ref{sec:longlived} \\
C18 & Positive-vote mirror, same specifications \emph{[AI, not yet reviewed]} & Requested at meeting 2 \\
\bottomrule
\end{longtable}
